% ---------------------------------------------------------------------------
% Appendix: the candidate sampler. Added 2026-11-05 (author: "we are not
% sampling randomly ... add this as a separate section in the paper"). The
% sampler was used in training and evaluation; earlier drafts said the
% candidates were drawn at random. Source: the author's specification,
% gaitnet-video/notes/sampler_coverage_author_2026-11-05.md. A numpy port of
% it (gaitnet-video/anim/gn_base.py, coverage_candidates) was checked for
% distinct safe cells, the pinned slot, determinism and mirror symmetry.
% ---------------------------------------------------------------------------
\section{Candidate Sampling}
\label{app:sampler}

At every tick, for each leg that may step, GaitNet places its
candidates so that they cover the leg's safe cells, instead of drawing
them at random. The sampler is deterministic and the same for every leg.

\paragraph{Objective.}
Let $S$ be the safe cells of the leg's grid
(Sec.~\ref{subsec:method-step}), in the yaw frame of the hip\verify{,
within the reach and box limits when these are set}. The sampler returns
$K = \corlTrainCandidatesPerLeg$ distinct safe cells, the candidates
$Q = \{\mathbf q_0, \dots, \mathbf q_{K-1}\}$, or all safe cells if
there are fewer, leaving the other slots empty. It aims at the quantization energy
\begin{equation}
E(Q) = \frac{1}{|S|} \sum_{\mathbf s \in S} \min_i
\|\mathbf s - \mathbf q_i\|^2,
\label{eq:sampler-energy}
\end{equation}
the $k$-means objective restricted to the safe cells, with the first
candidate pinned. A low energy means that, on average, a safe cell has a
candidate close by. The minimizer has no closed form, so the sampler
approximates it in the four steps below.

\paragraph{Nominal foothold.}
The first candidate is pinned near where a regular stride would put the
foot:
\begin{equation*}
\mathbf q_{\text{nom}} = T_{\text{nom}}\, \mathbf v_{\text{cmd}}
+ (0,\ y_{\text{foot}} - y_{\text{hip}}), \qquad
\mathbf q_0 = \arg\min_{\mathbf s \in S}
\|\mathbf s - \mathbf q_{\text{nom}}\|.
\end{equation*}
The first term follows the commanded base velocity $\mathbf v_{\text{cmd}}$, and the second keeps
the foot's current lateral distance from the hip. The horizon $T_{\text{nom}}$,
\corlNomHorizonS\,s, is a fixed prior, not a fitted value:
\corlNomLiftS\,s to lift-off, a \corlNomSwingS\,s swing, and half of a
stance as long as the swing, which puts the nominal foothold at
Raibert's neutral point at touchdown. The \corlNomLiftS\,s is part of
this prior only; steps still start at the tick they are chosen. The
first candidate does not move afterwards.

\paragraph{Spread.}
The other candidates are added one at a time by farthest-point sampling,
the greedy $k$-centre method of \citet{gonzalez_clustering_1985}: each
takes the safe cell farthest from all candidates chosen so far,
\begin{equation*}
\mathbf q_k = \arg\max_{\mathbf s \in S \setminus
\{\mathbf q_0, \dots, \mathbf q_{k-1}\}} \ \min_{i < k}
\|\mathbf s - \mathbf q_i\|, \qquad k = 1, \dots, K-1.
\end{equation*}
From any first point, the pinned one included, the covering radius
\begin{equation*}
r(Q) = \max_{\mathbf s \in S} \min_i \|\mathbf s - \mathbf q_i\|,
\end{equation*}
the largest distance from a safe cell to its nearest candidate, is then
at most twice the smallest radius that any $K$
cells reach \citep{gonzalez_clustering_1985}. Its picks, however, lie on
the boundary of the safe region, at the candidate margin of voids and at
the corners of the grid.

\paragraph{Refinement.}
\corlAnnealSteps{} steps of deterministic annealing
\citep{rose_deterministic_1998}, a soft $k$-means, then pull the
candidates inward. The temperature falls geometrically over the steps,
\begin{equation*}
\tau_t = \tau_0 \left(\tau_{\text{end}} / \tau_0\right)^{t/3},
\qquad t = 0, \dots, 3, \qquad
\tau_0 = (\corlAnnealTauStartCm\,\text{cm})^2, \quad
\tau_{\text{end}} = (\corlAnnealTauEndCm\,\text{cm})^2.
\end{equation*}
At step $t$, each safe cell $\mathbf s_j$ is shared among the candidates,
and each candidate moves to the mean of its share:
\begin{equation}
p_{ij} = \frac{\exp(-\|\mathbf s_j - \mathbf q_i\|^2 / \tau_t)}
{\sum_{k} \exp(-\|\mathbf s_j - \mathbf q_k\|^2 / \tau_t)}, \qquad
\tilde{\mathbf q}_i = \frac{\sum_j p_{ij}\, \mathbf s_j}{\sum_j p_{ij}}.
\label{eq:sampler-update}
\end{equation}
The pinned candidate keeps its cell, and a candidate that holds no share
stays where it is. Because the candidates compete for cells through the
softmax, no repulsion term is needed. The starting temperature is small
against the patch of cells each candidate covers, so the candidates do
not collapse onto one mean, the usual failure of annealing started too
hot.

\paragraph{Projection.}
After each step, every target returns to the nearest safe cell. When two
candidates claim one cell, the one added earlier keeps it and the other
takes the nearest cell not yet held. This greedy assignment, in the
order the candidates were added, approximates the closest distinct
assignment, ends after at most one round per candidate, and always finds
a cell, since there are at least as many safe cells as candidates.
Distances are computed in cell units, with targets rounded to
1/\corlAnnealRound{} of a cell, so they are exact, and ties go to the
lowest cell index. The candidates therefore do not depend on the batch
size or the device.

\paragraph{Left legs.}
For the two left legs, the sampler solves the leg's problem mirrored
left to right and maps the result back. Every step is
deterministic, so mirrored terrain and commands give mirrored
candidates.

\paragraph{Coverage.}
\verify{On grids with \corlCovVoidPct\% voids and 12.5\,cm tiles, the
covering radius was \corlCovRadiusCovCm\,cm with this sampler, against
\corlCovRadiusRandCm\,cm for the same number of cells drawn at random,
and the mean distance from a safe cell to its nearest candidate was
\corlCovMeanCovCm\,cm, against \corlCovMeanRandCm\,cm.}

% Ablation planned 2026-11-05 (~/dev/repos/gaitNET/ablation-sampler-plan.md):
% random-sampler retrains on the same seeds, same checkpoint rule, G125 protocol.
\paragraph{Ablation.}
\pending{Trained instead with the same number of cells drawn at random from the safe ones,
on the same xx training seeds and with the same checkpoint rule, GaitNet
succeeded in xx\% of the OCS2 benchmark's trials, against xx\% with
coverage for the same training seed on the same worlds, and in xx\%
against xx\% on the training controller in Isaac.}

\paragraph{What the refinement minimizes.}
Define the free energy
\begin{equation*}
F_\tau(Q) = -\tau \sum_{\mathbf s \in S} \log \sum_i
\exp\!\left(-\|\mathbf s - \mathbf q_i\|^2 / \tau\right).
\end{equation*}

\begin{proposition}[Annealing]
\label{prop:sampler-anneal}
(a)~For every set of candidates $Q$,
\begin{equation*}
|S|\, E(Q) - \tau |S| \log K \le F_\tau(Q) \le |S|\, E(Q).
\end{equation*}
(b)~With $m_i = \sum_j p_{ij}$, the update of
Eq.~\eqref{eq:sampler-update} is a scaled gradient step,
\begin{equation*}
\nabla_{\mathbf q_i} F_\tau = 2 m_i (\mathbf q_i - \tilde{\mathbf q}_i),
\qquad \text{so} \qquad
\tilde{\mathbf q}_i = \mathbf q_i - \frac{\nabla_{\mathbf q_i} F_\tau}{2 m_i}.
\end{equation*}
\end{proposition}

\begin{idmproof}{Proposition~\ref{prop:sampler-anneal}}
(a)~Fix $\mathbf s$, let $d_i = \|\mathbf s - \mathbf q_i\|^2$ and
$d = \min_i d_i$. Then $e^{-d/\tau} \le \sum_i e^{-d_i/\tau} \le K
e^{-d/\tau}$, so $d - \tau \log K \le -\tau \log \sum_i e^{-d_i/\tau}
\le d$. Summing over $S$ gives (a).
(b)~Differentiating, $\nabla_{\mathbf q_i} F_\tau = \sum_j p_{ij}\,
2(\mathbf q_i - \mathbf s_j) = 2 (m_i \mathbf q_i - \sum_j p_{ij}
\mathbf s_j) = 2 m_i (\mathbf q_i - \tilde{\mathbf q}_i)$.
\end{idmproof}

By (a), annealing minimizes a smoothed energy that sharpens into
Eq.~\eqref{eq:sampler-energy} as the temperature falls; at the last step
the gap is at most $\tau_{\text{end}} \log 16 \approx \corlAnnealGapCmSq$\,cm$^2$ per cell. By (b), each
update is a gradient step on the free energy, with a step size set by the candidate's share,
so there is no step size to tune, and at a fixed temperature it does not
increase the free energy, as for any expectation-maximization step
\citep{rose_deterministic_1998}. The projection after each step keeps
the candidates on safe cells but voids this guarantee.

\paragraph{Why coverage helps.}
Let $f$ score each safe cell as the policy would, and let
$\mathbf s^\star$ be the safe cell that $f$ scores highest.

\begin{proposition}[Cost of offering few cells]
\label{prop:sampler-regret}
If $f$ is $L$-Lipschitz on $S$, then
\begin{equation*}
f(\mathbf s^\star) - \max_{\mathbf q \in Q} f(\mathbf q) \le L\, r(Q).
\end{equation*}
If, in addition, $\mathbf s^\star$ is uniformly distributed on $S$, then
\begin{equation*}
\mathbb E\big[f(\mathbf s^\star) - \max_{\mathbf q \in Q} f(\mathbf
q)\big] \le L \sqrt{E(Q)}.
\end{equation*}
\end{proposition}

\begin{idmproof}{Proposition~\ref{prop:sampler-regret}}
Let $\mathbf q$ be the candidate nearest $\mathbf s^\star$. Then
$f(\mathbf s^\star) - \max_Q f \le f(\mathbf s^\star) - f(\mathbf q) \le
L \|\mathbf s^\star - \mathbf q\| \le L\, r(Q)$. For uniform
$\mathbf s^\star$, Jensen's inequality gives $\mathbb E[\min_{\mathbf q}
\|\mathbf s^\star - \mathbf q\|] \le (\mathbb E[\min_{\mathbf q}
\|\mathbf s^\star - \mathbf q\|^2])^{1/2} = \sqrt{E(Q)}$.
\end{idmproof}

Farthest-point sampling thus bounds the worst case to within a factor of
two, and the refinement lowers the bound on the average. The assumption
is strong, and we use the proposition only to motivate the sampler.
The score need not be Lipschitz where the terrain changes, at the edge of
a void. With the candidate mixer, it depends on the whole offered set.
And the bound concerns the offered choice, not where the controller lands
the foot.
